% Pursuing Stacks (cotes 134-2 to 134-8): archive cote 134-4, site page 72 (PDF page 73), typescript page 160.
% typescript. Draft compared: ps3.tex:3699-3759.
% AI-transcribed by Claude Opus 5.5 (claude-opus-5-5, adaptive thinking, effort high), Message Batches API, batch ps2 (msgbatch_01NVwr1Yuo8qJf5rykeGBXPK), 2026-11-12; not yet checked by eye.
\scanpage{134-4/72}
\mspage{160}
use of some list of conditions on $\underline{W}$ .

F\unsure{\hand{$\dagger$}}irst of all, I feel a review is needed of the \xout{terminologyxwexhavex}
few basic notions which have appeared in our work, relative to a set
of arrows $W\subset \mathrm{Fl}(M)$ in a general category M. We will not give to the
maps in W a specific name, such as ``weak equivalences'', as this may be
definitely misleading, in the general axiomatic set-up we want to deve-%
lop; \unsure{\deleted{w}here|where} $W_{(\mathrm{Cat})}$ is just one among many possible $\underline{W}$'s \deleted{the} and corres-%
pondingly for a small category A\unsure{\over{;}{,}}; $W_A = W_{A^{\wedge}}$ is just one among the many
$\underline{W}_A$'s , associated to the previous $\underline{W}$'s. When M = (Cat), it will be
understood we are working with a fixed set $\underline{W}\subset \mathrm{Fl}((\mathrm{Cat}))$, consisting of
the basic ``equivalences'', on which the whole modelizer story hinges.
\over{W}{W}e may call them $\underline{W}$-equivalences -- for the time being there will be no
question of varying $\underline{W}$ .

Coming back to a general pair (M, $W\subset \mathrm{Fl}(M)$) (not necessarily a
``modelizer''), we may call the maps in W W-equivalences. If M has a
final object e, we g\unsure{\over{o}{e}}t the corresponding notion of W-\uline{aspheric} \uline{object} of
M , namely an object such that the unique map
\[ X \longrightarrow e \]
is in W i.e. is a W-equivalence. We'll \inserted[typed]{define a} \xout{\unsure{say a map}} W-\uline{aspheric} \uline{map}
\[ f\colon X \longrightarrow Y \]
in M \hand{as} \deleted{W-aspheric if} one which is ``universally in W'' , by which I mean
that for any base-change
\[ Y' \longrightarrow Y \ , \]
the fiber product $X'=X\times_Y Y'$ exists (i.e. f is ``squarrable'') and the
map
\[ f'\colon X' \longrightarrow Y' \]
deduced from f by base change is in W. The thing to be quite careful
about is that for an object X in M, to say that X is W-\uline{aspheric over} e
(meaning that the \uline{map} $X \longrightarrow e$ is W-aspheric) implies X is \inserted{W-}aspheric \hand{i.e.}
\inserted{$X \to e$ is in W,} but the converse need not hold true. This causes a slight \inserted{psychological} \deleted{trouble} \hand{uneasy-}
\inserted{ness} \deleted{this terminology}, due to the fact I guess that the notion of an W-asphe-%
ric object has been defined after all in terms of the \uline{map} $X \longrightarrow e$ , and
consequently may be thought of as meaning \inserted{is} ``in W over e'' . Maybe we
should'nt use at all the word ``W-aspheric ob\unsure{\hand{j}}ect'' here, not even by
qualifying it as ``weakly W-aspheric'' to cause a feeling of caution, but
rather refer to this notion as `` X is a W-object'' -- and denote by \hand{M(W)}
the set of all those objects (or the corresponding full subcategory of M\hand{)},
and \xout{\unsure{say W}} call X $\underline{W}$-aspheric (dropping ``over e'')\hand{*} when it is ``universal-%
ly'' a W-object. The terminology we have been using so far was of \unsure{\over{c}{c}}ourse
\margin{* But we'll see immediately that this conflicts with the standard terminology in topoi -- so no good either!}
\margin{72}
